\section{Auswertung des Jahresabschlusses 2025}

Alle Betr\"age dieses Teils stammen aus dem Jahresbericht (Form 10-K) f\"ur
das am 31.\,Dezember 2025 abgelaufene Gesch\"aftsjahr. Jede Kennzahl wird
mit ihrer Bestimmungsgleichung eingef\"uhrt und anschlie{\ss}end
durchgerechnet, damit der Weg vom Ausgangswert zum Ergebnis nachpr\"ufbar
bleibt.

\subsection{Umsatz und Ver\"anderung zum Vorjahr}

Der Gesamtumsatz setzt sich aus Provisions- und Geb\"uhrenertr\"agen sowie
Zins- und sonstigen Ertr\"agen zusammen:
\begin{equation}
\label{eq:umsatz-zusammensetzung}
U = P + Z
\end{equation}
mit $U$ als Gesamtumsatz, $P$ als Provisions- und Geb\"uhrenertrag und $Z$
als Zins- und sonstigem Ertrag. F\"ur 2025 gilt
$\MioUSDexakt{\ProvisionsErtraege} + \MioUSDexakt{\InvestErtrag} =
\MioUSDexakt{\UmsatzGesamt}$ nach \MioUSDexakt{\UmsatzGesamtVJ} im
Vorjahr.\quelleGB{2025}{49}\quelleMerken{s5gb} Das entspricht einem Zuwachs
von \Pct{\UmsatzDelta}.

Diese Rate beschreibt jedoch nicht die Leistung des bestehenden
Gesch\"afts. Von der Zunahme der Kernprovisionen um
\MioUSDexakt{\OrganischNeugeschaeft} aus Neu- und Erneuerungsgesch\"aft
abgesehen stammt der gr\"o{\ss}te Teil aus Zuk\"aufen:
\MioUSDexakt{\UmsatzAusZukaeufenMDA} entfielen auf Gesellschaften, die im
Vorjahreszeitraum noch nicht zum Konzern geh\"orten; hinzu kamen
\MioUSDexakt{\WaehrungseffektMDA} aus der W\"ahrungsumrechnung, denen ein
Abgang ver\"au{\ss}erter Gesch\"afte von \MioUSDexakt{\DesinvestitionMDA}
gegen\"ubersteht.\quelleGB{2025}{36}\quelleMerken{s5gb}

\subsection{Organisches Umsatzwachstum als Fr\"uhindikator}
\label{sec:organisch}

Das Unternehmen ver\"offentlicht mit der \emph{Organic Revenue growth rate}
eine bereinigte Wachstumsrate: Kernprovisionen und Geb\"uhren ohne die in
den ersten zw\"olf Monaten erzielten Ertr\"age neu erworbener Einheiten,
ohne ver\"au{\ss}erte Gesch\"afte und ohne W\"ahrungseffekte.\quelleGB{2025}{35}\quelleMerken{s5gb}
Sie ist f\"ur einen Makler das, was f\"ur einen Anlagenbauer der
Auftragseingang ist: die Gr\"o{\ss}e, an der sich die Ertragskraft des
bestehenden Gesch\"afts ablesen l\"asst, bevor sie im Konzernumsatz sichtbar
wird.

Diese Rate ist 2025 von \Pct{\OrganischWachstumVJ} auf
\Pct{\OrganischWachstum} gefallen.\quelleGB{2025}{42}\quelleMerken{s5gb} Beide Segmente tragen
dazu bei, aber sehr ungleich: Im Retail-Segment ging sie von
\Pct{\RetailOrganischVJ} auf \Pct{\RetailOrganisch}
zur\"uck,\quelleGB{2025}{41}\quelleMerken{s5gb} in der Specialty Distribution von
\Pct{\SpezOrganischVJ} auf \Pct{\SpezOrganisch}.\quelleErneut{s5gb}
Der Einbruch liegt damit fast vollst\"andig im Spezialgesch\"aft. Als
Ursachen nennt das Unternehmen f\"ur den Einzelhandelsteil den Zeitpunkt
einmaliger Ertr\"age, nachlassende Pr\"amiensteigerungen, sinkende
Raten in einzelnen Deckungssparten und eine regulatorische \"Anderung im
britischen Gesch\"aft.\quelleErneut{s5gb}

\subsection{EBITDAC-Marge, berichtet und bereinigt}

Das Unternehmen f\"uhrt als operative Ergebnisgr\"o{\ss}e das EBITDAC, also
das Ergebnis vor Zinsen, Steuern, Abschreibungen und Ver\"anderung der
Erwerbspreisnachzahlungen. Die zugeh\"orige Marge ist
\begin{equation}
\label{eq:ebitdac-marge}
m = \frac{\mathit{EBITDAC}}{U}
\end{equation}
mit $m$ als Marge und $U$ als Gesamtumsatz nach
Gleichung~\eqref{eq:umsatz-zusammensetzung}. Berichtet ergibt sich
$\MioUSDexakt{\EBITDAC} / \MioUSDexakt{\UmsatzGesamt} =
\Pct{\EBITDACMarge}$ nach \Pct{\EBITDACMargeVJ} im
Vorjahr;\quelleGB{2025}{39}\quelleMerken{s6gb} bereinigt
\Pct{\EBITDACberMarge} nach \Pct{\EBITDACberMargeVJ}.\quelleGB{2025}{36}\quelleMerken{s6gb}

\needspace{5\baselineskip}
Das Verh\"altnis der beiden Gr\"o{\ss}en hat sich zwischen den Jahren
umgekehrt, und das ist der eigentliche Befund. 2024 lag das bereinigte
EBITDAC mit \MioUSDexakt{\EBITDACberVJ} \emph{unter} dem berichteten von
\MioUSDexakt{\EBITDACVJ}, weil die Bereinigung einen
Ver\"au{\ss}erungsgewinn von \MioUSDexakt{\VeraeusserungsergebnisVJ}
herausrechnete.\quelleGB{2025}{40}\quelleMerken{s6gb} 2025 liegt
es mit \MioUSDexakt{\EBITDACber} \emph{\"uber} dem berichteten von
\MioUSDexakt{\EBITDAC}, weil nun \MioUSDexakt{\AkqIntegrationskosten}
Akquisitions- und Integrationskosten hinzugerechnet und
\MioUSDexakt{\EscrowBewertung} aus der Bewertung der Treuhandverbindlichkeit
abgezogen werden.\quelleGB{2025}{51}\quelleMerken{s6gb}
\bk{Die bereinigte Marge steigt damit im selben Jahr, in dem die berichtete
f\"allt. Wer nur die bereinigte Reihe liest, sieht eine Verbesserung um
\Pct{\fpeval{round(\MargeDeltaBer,2)}}
Prozentpunkte, wer nur die berichtete liest, eine Verschlechterung um
\Pct{\fpeval{round(\MargeDeltaBerichtet,2)}} Prozentpunkte. Beide
Angaben sind richtig; die Differenz ist der Preis der \"Ubernahme.}

\subsection{Operatives EBITDAC}

Diese Analyse f\"uhrt im Weiteren das \emph{bereinigte} EBITDAC von
\MioUSDexakt{\EBITDACber} als operative Ergebnisgr\"o{\ss}e. Der Grund ist
die Vergleichbarkeit \"uber die Jahre: Die Akquisitions- und
Integrationskosten von \MioUSDexakt{\AkqIntegrationskosten} fallen einmalig
im Jahr des Erwerbs an, und die Bewertung der Treuhandverbindlichkeit folgt
dem eigenen Aktienkurs und nicht dem Gesch\"aft.\quelleErneut{s6gb}
Beide Posten sagen \"uber die Ertragskraft des laufenden Gesch\"afts nichts
aus. Der berichtete Wert von \MioUSDexakt{\EBITDAC} bleibt in
Tabelle~\ref{tab:kennzahlen} daneben stehen, damit die Bereinigung
sichtbar bleibt und nicht als Ergebnis erscheint.

Auf Segmentebene entfallen vom bereinigten EBITDAC
\MioUSDexakt{\RetailEBITDACber} auf Retail und
\MioUSDexakt{\SpezEBITDACber} auf die Specialty Distribution. Die Margen
liegen bei \Pct{\fpeval{round(\RetailMarge,1)}}
beziehungsweise
\Pct{\fpeval{round(\SpezMarge,1)}}: Das
Spezialgesch\"aft ist deutlich margenst\"arker als das
Einzelhandelsgesch\"aft.\quelleErneut{s6gb}

\subsection{Nettogewinn und Ver\"anderung zum Vorjahr}

Das den Anteilseignern zurechenbare Konzernergebnis stieg von
\MioUSDexakt{\KonzernergebnisVJ} auf \MioUSDexakt{\Konzernergebnis} und
damit um \Pct{\ErgebnisDelta}.\quelleGB{2025}{49}\quelleMerken{s6gb} Der Zuwachs bleibt weit
hinter dem Umsatzwachstum von \Pct{\UmsatzDelta} zur\"uck. Die Ursache
steht zwischen den beiden Gr\"o{\ss}en in der Gewinn- und Verlustrechnung:
Der Zinsaufwand stieg von \MioUSDexakt{\ZinsaufwandVJ} auf
\MioUSDexakt{\Zinsaufwand}, die Abschreibung auf erworbene immaterielle
Verm\"ogenswerte von \MioUSDexakt{\AbschrImmatVJ} auf
\MioUSDexakt{\AbschrImmat}.\quelleErneut{s6gb} Beide Zunahmen sind
unmittelbare Folgen der \"Ubernahme. Die Steuerquote betr\"agt
$\MioUSDexakt{\Ertragsteuern} / \MioUSDexakt{\ErgebnisVorSteuern} =
\Pct{\Steuerquote}$.

Je Aktie f\"allt das Ergebnis: verw\"assert \USDje{\EPSdil} nach
\USDje{\EPSdilVJ}, ein R\"uckgang um \PctBetrag{\EPSdilDelta}. Die Zahl der
umlaufenden Aktien ist 2025 um \num{\VerwaesserungStueck} Millionen St\"uck
gewachsen, also um \Pct{\AktienZuwachs}.\quelleGB{2025}{51}\quelleMerken{s7gb} Der h\"ohere
Gewinn verteilt sich damit auf so viel mehr Anteile, dass je Aktie weniger
\"ubrig bleibt als im Vorjahr.

\subsection{Dividende und Dividendenrendite}

Je Aktie wurden 2025 \USDje{\DividendeJeAktie} ausgesch\"uttet nach
\USDje{\DividendeJeAktieVJ} im Vorjahr, eine Erh\"ohung um
\Pct{\DividendeDelta}.\quelleWeb{Brown \& Brown, Inc., Dividend
History}{quelle:dividende}{28.08.2026} In der Kapitalflussrechnung
erscheinen daf\"ur \MioUSDexakt{\Dividendenzahlung} nach
\MioUSDexakt{\DividendenzahlungVJ}.\quelleGB{2025}{53}\quelleMerken{s7gb}
Die Dividendenrendite ist
\begin{equation}
\label{eq:dividendenrendite}
r_D = \frac{d}{K}
\end{equation}
mit $d$ als Dividende je Aktie und $K$ als Schlusskurs des Jahres. F\"ur
2025 ergibt sich $\USDje{\DividendeJeAktie} / \USDje{\KursSchlussZFV} =
\Pct{\Dividendenrendite}$ nach \Pct{\DividendenrenditeVJ} im Vorjahr. Die
Aussch\"uttungsquote auf das verw\"asserte Ergebnis je Aktie betr\"agt
\Pct{\Ausschuettungsquote}. F\"ur das erste Quartal 2026 hat das Board am
21.\,Januar 2026 eine Quartalsdividende von
\USDje{\DividendeQuartalNeu} beschlossen.\quelleGB{2025}{44}\quelleMerken{s7gb}

\subsection{Bilanz und Verschuldung}

Die Bilanzsumme ist 2025 von \MioUSDexakt{\BilanzsummeVJ} auf
\MioUSDexakt{\Bilanzsumme} gestiegen.\quelleErneut{s7gb} Der Zuwachs
liegt fast vollst\"andig in zwei Posten: Der Gesch\"afts- oder Firmenwert
w\"achst um \MioUSDexakt{\GoodwillZuwachs} auf \MioUSDexakt{\Goodwill}, die
erworbenen immateriellen Verm\"ogenswerte von
\MioUSDexakt{\ImmatVermoegenVJ} auf \MioUSDexakt{\ImmatVermoegen}. Der
Firmenwert allein macht damit \Pct{\GoodwillQuote} der Bilanzsumme aus und
\"ubersteigt das bilanzielle Eigenkapital um das
\Faktor{\GoodwillEKDeckung}-fache. Die Eigenkapitalquote betr\"agt
\Pct{\EKQuote}.

Die Nettoverschuldung ist definiert als
\begin{equation}
\label{eq:nettoverschuldung}
N = F - L
\end{equation}
mit $F$ als gesamten Finanzverbindlichkeiten und $L$ als eigenen
Zahlungsmitteln. Die Treuhandmittel von \MioUSDexakt{\TreuhandMittel}
bleiben dabei ausdr\"ucklich au{\ss}er Ansatz: Sie geh\"oren Kunden und
Versicherern und stehen dem Konzern zur Tilgung nicht zur
Verf\"ugung.\quelleGB{2025}{67}\quelleMerken{s7gb} Es gilt
$\MioUSDexakt{\Gesamtverschuldung} - \MioUSDexakt{\Zahlungsmittel} =
\MioUSDexakt{\Nettoverschuldung}$ nach
\MioUSDexakt{\NettoverschuldungVJ} im Vorjahr.\quelleErneut{s7gb}
Bezogen auf das bereinigte EBITDAC entspricht das dem
\Faktor{\Verschuldungsgrad}-fachen nach \Faktor{\VerschuldungsgradVJ} im
Vorjahr, wobei EBITDAC das operative Ergebnis \emph{vor} Abschreibungen
bezeichnet und nicht den Gewinn.

\bk{Dieser Sprung ist die zentrale bilanzielle Ver\"anderung des Jahres.
Er ist jedoch nur zur H\"alfte aussagekr\"aftig, weil der Nenner das
Ergebnis eines Jahres ist, in dem die erworbene Gesellschaft erst ab
1.\,August konsolidiert wurde. Auf das bereinigte EBITDAC der
Pro-forma-Erl\"ose von \MioUSD{\EBITDACProForma} bezogen liegt der
Verschuldungsgrad bei \Faktor{\VerschuldungsgradProForma}. Teil~\ref{sec:flow}
rechnet mit dieser zweiten Gr\"o{\ss}e weiter und begr\"undet das dort.}

\subsection*{\"Uberblick der Kennzahlen 2024 und 2025}
\addcontentsline{toc}{subsection}{\"Uberblick der Kennzahlen 2024 und 2025}

\begin{table}[htbp]
\centering
\caption{Kennzahlen im Zweijahresvergleich. Betr\"age in Mio.~USD, sofern
nicht anders angegeben.}
\label{tab:kennzahlen}
\begin{tabular}{lrrr}
\toprule
Kennzahl & 2024 & 2025 & Ver\"anderung\\
\midrule
Gesamtumsatz                  & \num{\UmsatzGesamtVJ}    & \num{\UmsatzGesamt}    & \Pct{\fpeval{round(\UmsatzDelta,1)}}\\
Kernprovisionen               & \num{\KernprovisionenVJ} & \num{\Kernprovisionen} & \Pct{\fpeval{round(\KernprovisionenDelta,1)}}\\
Organisches Wachstum in \%    & \num{\OrganischWachstumVJ} & \num{\OrganischWachstum} & --\\
EBITDAC berichtet             & \num{\EBITDACVJ}         & \num{\EBITDAC}         & \Pct{\fpeval{round(\EBITDACDelta,1)}}\\
EBITDAC bereinigt             & \num{\EBITDACberVJ}      & \num{\EBITDACber}      & \Pct{\fpeval{round(\EBITDACberDelta,1)}}\\
EBITDAC-Marge bereinigt in \% & \num{\fpeval{round(\EBITDACberMargeVJ,1)}} & \num{\fpeval{round(\EBITDACberMarge,1)}} & --\\
Konzernergebnis               & \num{\KonzernergebnisVJ} & \num{\Konzernergebnis} & \Pct{\fpeval{round(\ErgebnisDelta,1)}}\\
Ergebnis je Aktie in USD      & \num{\EPSdilVJ}          & \num{\EPSdil}          & \Pct{\fpeval{round(\EPSdilDelta,1)}}\\
Dividende je Aktie in USD     & \num{\DividendeJeAktieVJ} & \num{\DividendeJeAktie} & \Pct{\fpeval{round(\DividendeDelta,1)}}\\
Operativer Cash-Flow          & \num{\OperativerCFVJ}    & \num{\OperativerCF}    & \Pct{\fpeval{round(\OperativerCFDelta,1)}}\\
Nettoverschuldung             & \num{\NettoverschuldungVJ} & \num{\Nettoverschuldung} & --\\
Nettoverschuldung/EBITDAC     & \num{\fpeval{round(\VerschuldungsgradVJ,2)}} & \num{\fpeval{round(\Verschuldungsgrad,2)}} & --\\
Eigenkapital                  & \num{\EigenkapitalVJ}    & \num{\Eigenkapital}    & \Pct{\fpeval{round(\EigenkapitalZuwachs,1)}}\\
Schlusskurs in USD            & \num{\KursSchlussZFIV}   & \num{\KursSchlussZFV} & \Pct{\fpeval{round(\KursDeltaZweiFuenf,1)}}\\
\bottomrule
\end{tabular}
\end{table}

\FloatBarrier
